\documentclass{article}
\usepackage{amsmath}
\usepackage{amssymb}
\usepackage{geometry}
\usepackage{graphicx}
\usepackage{xepersian}
\settextfont{XB Niloofar}

\geometry{a4paper, margin=1in}

\title{\lr{The Kernel Method (Kernel Trick)} \\ (روش هسته / ترفند هسته)}
\author{}
\date{}

\begin{document}

\maketitle

\section*{فهرست مطالب}
\begin{enumerate}
    \item مقدمه
    \item محدودیت‌های مدل‌های خطی
    \item نگاشت ویژگی (\lr{Feature Mapping})
    \item ترفند هسته (\lr{The Kernel Trick})
    \item الگوریتم‌های استفاده‌کننده از هسته
    \item توابع هسته رایج (\lr{Common Kernel Functions})
    \item قضیه مرسر (\lr{Mercer's Theorem})
    \item تفسیرهای بیشتر و عملیات در فضای ویژگی
    \item مزایا و معایب ترفند هسته
    \item خوشه‌بندی کی-مینز هسته (\lr{Kernel K-means Clustering})
\end{enumerate}

\newpage

\section*{۱. مقدمه}
بسیاری از الگوریتم‌های یادگیری ماشین از جمله رگرسیون خطی، رگرسیون لجستیک، تحلیل مولفه‌های اصلی (\lr{PCA}) و دسته‌بندی‌کننده‌های خطی (مانند پرسپترون یا ماشین‌های بردار پشتیبان خطی) بر روابط خطی بین نقاط داده تکیه دارند. این الگوریتم‌ها زمانی که داده‌ها به صورت خطی تفکیک‌پذیر باشند یا بتوان آنها را به طور مناسب با توابع خطی مدل‌سازی کرد، عملکرد خوبی دارند. با این حال، بسیاری از مسائل دنیای واقعی شامل ساختارهای داده غیرخطی و پیچیده‌ای هستند که مدل‌های خطی نمی‌توانند آنها را به طور موثر ثبت کنند.
\begin{figure}[h!]
  \centering
  \includegraphics[width=0.75\textwidth]{img_page_1_1.png}
  \caption{مفهوم هندسی نگاشت داده‌ها به فضای با بعد بالا در روش‌های هسته}
\end{figure}




روش هسته (\lr{Kernel Method}) یا ترفند هسته (\lr{Kernel Trick}) راه قدرتمندی برای اعمال الگوریتم‌های خطی به داده‌های غیرخطی ارائه می‌دهد. ایده اصلی نگاشت داده‌ها به یک فضای ویژگی با ابعاد بالاتر ($\mathcal{F}$) است که در آن داده‌ها به صورت خطی تفکیک‌پذیر یا برای مدل‌سازی خطی آسان‌تر می‌شوند. ترفند این است که این نگاشت می‌تواند به طور ضمنی با استفاده از یک تابع هسته انجام شود، بدون اینکه نیازی به محاسبه صریح مختصات نقاط داده در فضای ویژگی با ابعاد بالا (و به طور بالقوه نامتناهی) باشد. این کار از بار محاسباتی مرتبط با نگاشت صریح جلوگیری می‌کند.

\section*{۲. محدودیت‌های مدل‌های خطی}
در یک مسئله دسته‌بندی باینری، یک مدل خطی تلاش می‌کند تا ابرصفحه‌ای برای جداسازی نقاط داده متعلق به کلاس‌های مختلف پیدا کند. در فضای ۲بعدی یک خط ($w^T x + b = 0$)، در ۳بعدی یک صفحه ($w_1 x_1 + w_2 x_2 + w_3 x_3 + b = 0$) و در فضای $p$-بعدی یک زیرفضای $(p-1)$-بعدی است. اگر چنین ابرصفحه جداکننده‌ای وجود داشته باشد، داده‌ها به صورت خطی تفکیک‌پذیر نامیده می‌شوند. اما اگر داده‌ها ساختار غیرخطی داشته باشند (مانند مسئله XOR یا دایره‌های هم‌مرکز)، هیچ ابرصفحه واحدی نمی‌تواند کلاس‌ها را به درستی در فضای ورودی اصلی جدا کند.

\section*{۳. نگاشت ویژگی}
برای مدیریت داده‌های غیرخطی، می‌توان داده‌ها را به یک فضای ویژگی جدید و با ابعاد بالاتر منتقل کرد که در آنجا ممکن است به صورت خطی تفکیک‌پذیر شوند. این کار با استفاده از یک تابع نگاشت که با $\phi$ نشان داده می‌شود، انجام می‌گیرد:
\[
\phi: \mathcal{X} \to \mathcal{F}
\]
که در آن $\mathcal{X}$ فضای ورودی اصلی (به عنوان مثال $\mathbb{R}^p$) و $\mathcal{F}$ فضای ویژگی با ابعاد بسیار بالاتر (به عنوان مثال $\mathbb{R}^d$ که $d \gg p$ است یا حتی نامتناهی) است.

ترفند هسته با مفهوم هسته‌سازی (\lr{kernelization}) در طراحی الگوریتم‌های کارآمد (که در آن ورودی‌ها با ورودی کوچکتری به نام هسته جایگزین می‌شوند) متفاوت است.

\section*{۴. ترفند هسته}
قدرت روش هسته در این مشاهدات نهفته است که بسیاری از الگوریتم‌های خطی را می‌توان به گونه‌ای فرموله کرد که فقط به ضرب داخلی بین نقاط داده نیاز داشته باشند ($x_i^T x_j$).
اگر از نگاشت ویژگی $\phi$ استفاده کنیم، این الگوریتم‌ها باید ضرب داخلی را در فضای ویژگی با ابعاد بالا محاسبه کنند: $\phi(x_i)^T \phi(x_j)$. محاسبه صریح $\phi(x)$ به ویژه اگر ابعاد آن نامتناهی یا بسیار بزرگ باشد، می‌تواند از نظر محاسباتی بسیار گران تمام شود.

ترفند هسته از این محاسبات صریح جلوگیری می‌کند. ما یک تابع هسته تعریف می‌کنیم که ضرب داخلی در فضای ویژگی را مستقیماً از روی بردارهای ورودی اصلی محاسبه کند:
\[
K(x, z) = \phi(x)^T \phi(z)
\]
اگر بتوانیم تابع $K$ را پیدا کنیم که این ضرب داخلی را بدون محاسبه اولیه $\phi(x)$ و $\phi(z)$ به طور کارآمد محاسبه کند، به سادگی هر ضرب داخلی را با ارزیابی تابع هسته جایگزین می‌کنیم. این به ما امکان می‌دهد در فضای ویژگی با ابعاد بالا عمل کنیم در حالی که محاسبات را فقط در فضای ورودی اصلی انجام می‌دهیم.

ضرب داخلی نشان‌دهنده شباهت بین نقاط در فضای ویژگی است. فضای ویژگی $\mathcal{F}$ به عنوان فضای هیلبرت هسته بازتولید‌کننده (\lr{Reproducing Kernel Hilbert Space} یا \lr{RKHS}) شناخته می‌شود.

\section*{۵. الگوریتم‌های استفاده‌کننده از هسته}
هر الگوریتمی که بتواند صرفاً بر حسب ضرب داخلی بین نمونه‌های ورودی بیان شود، پتانسیل هسته‌دار شدن (\lr{kernelized}) را دارد: دسته‌بندی‌کننده‌های خطی (پرسپترون، رگرسیون لجستیک)، تحلیل مولفه‌های اصلی هسته (\lr{Kernel PCA})، ماشین‌های بردار پشتیبان (\lr{SVM})، فرآیندهای گاوسی، رگرسیون ستیغی هسته و روش‌های خوشه‌بندی.
\begin{figure}[h!]
  \centering
  \includegraphics[width=0.75\textwidth]{img_page_4_1.png}
  \caption{نمایش نگاشت ویژگی غیرخطی و محاسبات ضرب داخلی ساده با استفاده از تابع هسته}
\end{figure}




\section*{۶. توابع هسته رایج}
\begin{itemize}
    \item \textbf{هسته خطی (\lr{Linear Kernel}):} 
    \[
    K(x, z) = x^T z
    \]
    با نگاشت همانی $\phi(x)=x$ که الگوریتم خطی اصلی را بازیابی می‌کند.
    \item \textbf{هسته چندجمله‌ای (\lr{Polynomial Kernel}):} 
    \[
    K(x, z) = (\gamma x^T z + c)^d
    \]
    که در آن $\gamma، c \ge 0$ و درجه چندجمله‌ای $d$ ابرپارامتر هستند. این هسته متناظر با فضای ویژگی شامل ترکیب‌های چندجمله‌ای ویژگی‌های اصلی تا درجه $d$ است.
    \item \textbf{هسته گاوسی / تابع پایه شعاعی (\lr{Gaussian Kernel / RBF}):} 
    \[
    K(x, z) = \exp\left( -\frac{\|x-z\|^2}{2\sigma^2} \right) = \exp(-\gamma \|x-z\|^2)
    \]
    که در آن $\gamma = 1/(2\sigma^2)$ یک ابرپارامتر است. این هسته متناظر با فضای ویژگی با ابعاد نامتناهی است.
\end{itemize}

\section*{۷. قضیه مرسر: چه زمانی \lr{K} یک هسته معتبر است؟}
هر تابعی نمی‌تواند به عنوان ضرب داخلی در یک فضای هیلبرت تفسیر شود. قضیه مرسر شرایط لازم برای معتبر بودن یک تابع هسته را فراهم می‌کند.

\subsection*{ماتریس گرام (\lr{Gram Matrix})}
با داشتن مجموعه‌داده $\{x_1, \dots, x_n\}$، ماتریس گرام $K$ یک ماتریس $n \times n$ است که در آن عنصر $K_{ij}$ ضرب داخلی بردارهای ویژگی است: $K_{ij} = \langle \phi(x_i), \phi(x_j) \rangle = K(x_i, x_j)$.

\subsection*{ماتریس معین مثبت (\lr{Positive Semidefinite})}
یک ماتریس متقارن $K$ معین مثبت است اگر برای هر بردار غیرصفر $c \in \mathbb{R}^n$ داشته باشیم: $c^T K c \ge 0$.

\subsection*{قضیه مرسر (ساده‌شده)}
یک تابع متقارن و پیوسته $K$ یک هسته معتبر است اگر و تنها اگر ماتریس گرام $K$ برای هر مجموعه متناهی از نقاط، معین مثبت باشد.

\textbf{اثبات یک جهت قضیه (اگر $K$ هسته باشد، ماتریس گرام PSD است):}
فرض کنید $K(x, z) = \phi(x)^T \phi(z)$. برای هر بردار $c = (c_1, \dots, c_n)^T \in \mathbb{R}^n$ داریم:
\[
c^T K c = \sum_{i=1}^n \sum_{j=1}^n c_i c_j K_{ij} = \sum_{i=1}^n \sum_{j=1}^n c_i c_j (\phi(x_i)^T \phi(x_j))
\]
\[
c^T K c = \left( \sum_{i=1}^n c_i \phi(x_i) \right)^T \left( \sum_{j=1}^n c_j \phi(x_j) \right) = \|v\|^2 \ge 0
\]
که در آن $v = \sum c_i \phi(x_i)$ برداری در فضای ویژگی $\mathcal{F}$ است. چون هنجار مربع هر بردار همواره غیرمنفی است، پس ماتریس گرام معین مثبت است.

\section*{۸. تفسیرهای بیشتر و عملیات در فضای ویژگی}
\begin{itemize}
    \item \textbf{هنجار (نرم) یک نقطه در فضای ویژگی:} 
    \[
    \|\phi(x)\|^2 = \phi(x)^T \phi(x) = K(x, x)
    \]
    \item \textbf{میانگین در فضای ویژگی:} برای داده‌های $\{x_1, \dots, x_n\}$ میانگین در فضای ویژگی $m_{\phi} = \frac{1}{n} \sum \phi(x_i)$ است. ضرب داخلی آن با هر نقطه: $\phi(x)^T m_{\phi} = \frac{1}{n} \sum K(x, x_j)$.
    \item \textbf{فاصله بین نقاط در فضای ویژگی:} 
    \[
    \|\phi(x_i) - \phi(x_j)\| = \sqrt{K(x_i, x_i) + K(x_j, x_j) - 2 K(x_i, x_j)}
    \]
    \item \textbf{واریانس کل در فضای ویژگی:}
    \[
    Var_{\phi} = \frac{1}{n} \sum_{i=1}^n \|\phi(x_i) - m_{\phi}\|^2 = \frac{1}{n} \sum_{i=1}^n K(x_i, x_i) - \frac{2}{n^2} \sum_{i,j} K(x_i, x_j) + \frac{1}{n^2} \sum_{k,l} K(x_k, x_l)
    \]
    \item \textbf{مرکزی‌سازی (\lr{Centering}):} با تعریف $\tilde{\phi}(x_i) = \phi(x_i) - m_{\phi}$، ماتریس هسته مرکزی‌شده $\tilde{K}$ از رابطه زیر به دست می‌آید که گام اصلی در Kernel PCA است:
    \[
    \tilde{K} = K - \mathbf{1}_n K - K \mathbf{1}_n + \mathbf{1}_n K \mathbf{1}_n
    \]
    که در آن $\mathbf{1}_n$ ماتریس با مقادیر $1/n$ است.
    \item \textbf{نرمال‌سازی هنجار واحد:} هسته نرمال‌سازی شده (یا هسته کسینوسی) به صورت زیر است:
    \[
    K_{\text{norm}}(x, z) = \frac{K(x, z)}{\sqrt{K(x, x)} \sqrt{K(z, z)}}
    \]
\end{itemize}

\section*{۹. مزایا و معایب ترفند هسته}
\begin{itemize}
    \item \textbf{مزایا:} توانایی مدل‌سازی غیرخطی بودن، کارایی محاسباتی با عدم محاسبه صریح $\phi(x)$، ماژولار بودن و امکان کار با داده‌های غیربرداری (مانند متون یا گراف‌ها).
    \item \textbf{معایب:} هزینه محاسباتی در داده‌های بزرگ (محاسبه ماتریس گرام زمان $O(n^2 d)$ و فضای $O(n^2)$ می‌برد)، دشواری انتخاب هسته و ابرپارامترها و کاهش تفسیرپذیری مرز تصمیم‌گیری.
\end{itemize}

\section*{۱۰. خوشه‌بندی کی-مینز هسته}
الگوریتم \lr{K-means} استاندارد فقط خوشه‌های با مرزهای خطی را پیدا می‌کند. برای خوشه‌های با ساختار غیرخطی، کی-مینز هسته (\lr{Kernel K-means}) با ترفند هسته داده‌ها را به فضای ویژگی با ابعاد بالاتر نگاشت می‌کند.

فاصله مربع بین یک نقطه و مرکز خوشه $C_i$ در فضای ویژگی به صورت زیر محاسبه می‌شود:
\[
\|\phi(x_j) - \mu^{\phi}_i\|^2 = K(x_j, x_j) - \frac{2}{n_i} \sum_{x_a \in C_i} K(x_a, x_j) + \frac{1}{n_i^2} \sum_{x_a, x_b \in C_i} K(x_a, x_b)
\]
در این الگوریتم، پس از تخصیص تصادفی اولیه نقاط به خوشه‌ها، فواصل در فضای ویژگی با استفاده از رابطه فوق محاسبه شده و نقاط به نزدیک‌ترین خوشه تخصیص داده می‌شوند. این فرآیند تا همگرایی تکرار می‌شود.

\end{document}
